% Part: many-valued-logic
% Chapter: sequent-calculus
% Section: rules-and-proofs

\documentclass[../../../include/open-logic-section]{subfiles}

\begin{document}

\olfileid{mvl}{seq}{rul}

\olsection{નિયમો અને \usetoken{P}{derivation}}

નીચેના માટે, $\Gamma, \Delta, \Pi, \Lambda$ એ !!{sentence}sની
સાન્ત શ્રેણીઓ દર્શાવે તેમ માનો.

\begin{defn}[સિક્વન્ટ]
\emph{$n$-બાજુવાળો સિક્વન્ટ} નીચેના સ્વરૂપનું પદ છે:
\[
\Gamma_1 \nSequent \dots \nSequent \Gamma_n
\]
અહીં દરેક $\Gamma_i$ એ ભાષા~$\Lang L$નાં !!{sentence}sની
સાન્ત (સંભવતઃ ખાલી) શ્રેણી છે.
\end{defn}
\sourcecorrection{OLMV-019}{સ્થિર અંગ્રેજી મૂળ
સિક્વન્ટનાં $n$ સ્થાનો દર્શાવ્યા પછી દરેક સ્થાનને
$\Gamma_1$ કહે છે. સામાન્ય સ્થાન માટે $\Gamma_i$
વાપરીને તેની વ્યાખ્યા અને ગણિતીય સૂચકાંક એકરૂપ કર્યા છે.}

\begin{defn}[આરંભિક સિક્વન્ટ]
\emph{$n$-બાજુવાળો આરંભિક સિક્વન્ટ} ભાષાના કોઈપણ
!!{sentence} $!A$ માટે $!A \nSequent \dots \nSequent !A$
સ્વરૂપનો $n$-બાજુવાળો સિક્વન્ટ છે.

જો ભાષામાં $0$-સ્થાનવાળો સંયોજક~$\star$ હોય, એટલે
કે વિધાનાત્મક અચલ હોય, તો $\dots \nSequent
\star \nSequent \dots$ સિક્વન્ટ પણ આરંભિક ગણીએ છીએ.
તેમાં $\star$ તેના $\tf{\star} \in V$ સાથે સંકળાયેલા
સત્યમૂલ્યના સ્થાનમાં આવે છે અને અન્ય સ્થાનો ખાલી રહે છે.
\end{defn}

$n$-મૂલ્ય તર્કશાસ્ત્ર~$\Log{L}$ના દરેક સંયોજક માટે,
તે સંયોજક~$\Log{L}$માં લઈ શકે એવા દરેક સત્યમૂલ્યને
અનુરૂપ એક તાર્કિક નિયમ હોય છે. $\Log{L}$ના $n$-બાજુવાળા
સિક્વન્ટ કલનમાં !!^{derivation}s એ સિક્વન્ટોનાં વૃક્ષો છે:
સૌથી ઉપરના સિક્વન્ટો આરંભિક સિક્વન્ટો હોય છે, અને
જો કોઈ સિક્વન્ટ બીજા એક કે વધુ સિક્વન્ટોની નીચે હોય,
તો તે $\Log{L}$ના સંયોજકોના અનુમાન-નિયમ વડે તેમાંથી
યોગ્ય રીતે મળવું જોઈએ.

\begin{defn}[પ્રમેયો]
!!^a{sentence}~$!A$ એ $n$-મૂલ્ય તર્કશાસ્ત્ર~$\Log{L}$નું
\emph{પ્રમેય} છે જો $\Log{L}$ના નિર્દિષ્ટ સત્યમૂલ્યને
અનુરૂપ દરેક સ્થાનમાં $!A$ ધરાવતા $n$-સિક્વન્ટની
!!a{derivation} હોય. જો $!A$ પ્રમેય હોય તો
$\Proves[\Log{L}] !A$ લખીએ છીએ, અને ન હોય તો
$\Proves/[\Log{L}] !A$ લખીએ છીએ.
\end{defn}

\begin{defn}[!!^{derivability}]
!!^a{sentence}~$!A$ એ !!{sentence}sના ગણ~$\Gamma$માંથી
$n$-મૂલ્ય તર્કશાસ્ત્ર~$\Log{L}$માં \emph{!!{derivable}}
છે, $\Gamma \Proves[\Log{L}] !A$, જો અને માત્ર જો
કોઈ સાન્ત ઉપગણ~$\Gamma_0 \subseteq \Gamma$
અને $\Gamma_0$નાં !!{sentence}sની કોઈ શ્રેણી $\Gamma_0'$
હોય જેથી નીચેના સિક્વન્ટની !!a{derivation} હોય:
\[ \Lambda_1 \nSequent \dots \nSequent \Lambda_n \]
અહીં જો સ્થાન $i$ કોઈ નિર્દિષ્ટ સત્યમૂલ્યને અનુરૂપ હોય
તો $\Lambda_i$ એ $!A$ છે, અને અન્યથા $\Gamma_0'$ છે.
જો $!A$ એ~$\Gamma$માંથી !!{derivable} ન હોય તો
$\Gamma \Proves/ !A$ લખીએ છીએ.
\end{defn}

દાખલા તરીકે, $3$-મૂલ્ય લુકાશેવિચ તર્કશાસ્ત્રમાં
$3$-બાજુવાળું સિક્વન્ટ કલન છે. $3$-બાજુવાળા સિક્વન્ટ
$\Gamma \nSequent \Pi \nSequent \Delta$માં $\Gamma$
એ~$\False$ને, $\Delta$ એ~$\True$ને, અને $\Pi$
એ~$\Undef$ને અનુરૂપ છે. સ્વયંસિદ્ધિઓ
$!A \nSequent !A \nSequent !A$ છે. માત્ર $\True$
નિર્દિષ્ટ હોવાથી, $\Gamma \Proves[\LogLuk[3]] !A$
ત્યારે અને માત્ર ત્યારે જ જ્યારે
$\Gamma \nSequent \Gamma \nSequent !A$ સિક્વન્ટની
!!a{derivation} હોય. (જો $\Undef$ પણ નિર્દિષ્ટ હોત,
તો $\Gamma \nSequent !A \nSequent !A$ની
!!a{derivation} જોઈતી હોત.)

\end{document}
